\section*{Sammendrag}
\addcontentsline{toc}{section}{Sammendrag}
Selv om kunstig intelligens har oppnådd bemerkelsesverdige resultater det siste tiåret, står feltet fortsatt overfor flere grunnleggende utfordringer. KI krever store mengder energi, har vanskeligheter med å tilegne seg ny kunnskap uten å overskrive det som allerede er lært, generaliserer ofte dårlig til nye situasjoner som skiller seg fra treningsdataene, og mangler den kausale forståelsen av verden som utvikles gjennom erfaring og samhandling med den fysiske verden. Dersom vi ønsker å løse disse problemene, er det naturlig å hente inspirasjon fra biologien, og særlig fra den menneskelige hjernen, hvor mange av disse egenskapene allerede finnes. Hvis fremtidige KI-systemer i større grad skal bygge på biologisk inspirerte arkitekturer, trenger de også læringsalgoritmer som kan fungere innenfor slike arkitekturer. Backpropagation trener kunstige nevrale nettverk svært effektivt, men er vanskelig å forene med biologisk læring. Dette motiverer søket etter biologisk plausible læringsalgoritmer.

Denne masteroppgaven studerer perturbasjonsbasert læring, en retning innen biologisk plausibel læring der tilfeldig støy legges til et nettverks vekter eller nevrale aktivitet. Endringen i ytelse som følger av denne støyen, brukes deretter til å avgjøre om perturbasjonene skal forsterkes eller reverseres. Oppgaven gjennomgår først den eksisterende litteraturen innenfor dette feltet, hvor den identifiserer et uavklart spørsmål knyttet til forholdet mellom de to hovedformene for perturbasjonsbasert læring. Med utgangspunkt i en idé fra bayesianske nevrale nettverk adresseres dette spørsmålet deretter gjennom utledningen av et teorem som viser at den ene familien av læringsmetoder kan forstås som en variansredusert versjon av den andre. I denne utledningen introduseres også en ny metode kalt input-skalert nodeperturbasjon. Så vidt vi kjenner til, har verken det teoretiske resultatet eller metoden blitt utledet tidligere.

De empiriske resultatene gir støtte for at teoremet også har praktisk betydning. Etter hvert som oppgavene blir mer komplekse, viser metodene med lavere varians tydelig bedre læringsytelse enn den andre formen. Resultatene viser også at input-skalert nodeperturbasjon fungerer bedre enn de eksisterende perturbasjonsmetodene den sammenlignes med, særlig på de mer utfordrende oppgavene. Dette tyder på at det teoretiske bidraget også fører til en metode med forbedret læringsytelse. Med det sagt fungerer backpropagation fortsatt langt bedre enn samtlige perturbasjonsbaserte algoritmer på de mest utfordrende oppgavene. Denne masteroppgaven fremstiller derfor ikke perturbasjonsbasert læring som en erstatning for backpropagation, men som en biologisk motivert læringsalgoritme som er kapabel til å lære ikke-trivielle oppgaver, samtidig som den har tydelig rom for videre utvikling, inkludert den typen teoretisk motivert forbedring som utvikles i dette arbeidet.

\clearpage
